\chapter{高等技術}
\label{ch:extendedtechniques}

到目前為止，我們用的型別大致上是 \texttt{int}、\texttt{char}、\texttt{double}
這些基本型別，加上陣列與函式。 實務上還有幾種 C 語言提供的「進階型別
工具」。 \texttt{typedef}、\texttt{enum}、\texttt{union}。 以及跟「變數能活
多久、誰能看到、在記憶體裡怎麼排」有關的 storage class 與對齊規則。 這些
觀念在題庫裡常常跟位元操作一起出現（例如用 \texttt{typedef} 包一個集合型
別、再用 \texttt{unsigned long long} 存 bitset），但本身並不是位元技巧；
這一章先把型別與記憶體布局的基礎打好，下一章再專心處理位元運算。

\section{typedef：替型別取別名}
\label{sec:ext-typedef}

\texttt{typedef} 讓我們替既有的型別取一個新名字。 編譯器看見
\texttt{Score} 時，仍然把它當成 \texttt{int} 來處理。 不會多占記憶體、也不
會改變運算規則；差別只在於程式碼讀起來更像在描述「這個 \texttt{int} 代表
什麼意義」。 例如考試分數、座標、計數器，都可以各自 typedef 成不同的名
字，讓函式簽名一眼看出參數的語意。

\begin{code}{c}{替 int 取別名 Score}
typedef int Score;

Score a, b;
scanf("%d %d", &a, &b);
printf("%d\n", a + b);
\end{code}

範例 \thechapter.2 把兩個分數讀進來相加。 如果之後題目改成「平均分要印
小數」，我們只需要把 \texttt{typedef int Score} 改成
\texttt{typedef double Score}，其餘變數名稱 \texttt{a}、\texttt{b} 不必
全部重命名。 這正是 typedef 在維護大型程式時的價值：型別的「物理表示」可以
換，但「語意上的名字」可以留著。

\begin{example}{c}{(type-alias.c) 用 Score 表示整數分數}
#include <stdio.h>

typedef int Score;

int main(void) {
    Score a, b;
    scanf("%d %d", &a, &b);
    printf("%d\n", a + b);
}
\end{example}

\iosample{
3 4 \\
}{
7 \\
}

\keypoint{\texttt{typedef} 不創造新型別，只是替既有型別取別名；編譯器仍
依底層型別決定大小與運算，但程式碼裡的名字可以表達設計意圖。}

題庫裡較後段的集合 API 會用 \texttt{typedef struct set \{ \ldots \} Set;}
把「內部怎麼存 bitset」藏起來，對外只暴露 \texttt{Set} 這個名字。  typedef
在這裡扮演的角色，跟本節的 \texttt{Score} 完全一樣，只是包裝的底層型別
從 \texttt{int} 換成了 \texttt{struct}。

\begin{drills}[label=drills:ext-typedef]
\item \texttt{typedef int Score;} 之後寫 \texttt{Score a;}，\texttt{a}
實際的底層型別是什麼？ 編譯器看到 \texttt{Score} 時怎麼處理它？
\item 若要把 \texttt{Score} 底層型別從 \texttt{int} 換成 \texttt{double}，
需要改幾處程式碼？ 這是 \texttt{typedef} 在維護程式時的什麼價值？
\item 題庫裡 \texttt{typedef struct set \{ ... \} Set;} 跟本節的
\texttt{typedef int Score;}，兩者裡 \texttt{typedef} 扮演的角色一樣嗎？
\end{drills}

\section{enum：具名整數常數}
\label{sec:ext-enum}

有時候程式需要在少數幾個固定選項之間做選擇。 例如兩位候選人誰得票較
多、星期幾、方向上下左右。 我們當然可以直接用 \texttt{0}、\texttt{1}、
\texttt{2} 這些魔術數字，但讀程式的人得靠註解或變數名才能猜出
\texttt{1} 到底代表什麼。 \texttt{enum}（enumeration，列舉）讓編譯器自
動替這些選項分配整數常數，並給每個常數一個名字。

\begin{code}{c}{宣告列舉型別與列舉常數}
typedef enum { CAND_A, CAND_B } Candidate;

Candidate winner;
if (votesA > votesB) {
    winner = CAND_A;
} else {
    winner = CAND_B;
}
\end{code}

沒有特別指定值的時候，第一個列舉常數是 \texttt{0}，之後依序遞增：
\texttt{CAND\_A} 是 \texttt{0}、\texttt{CAND\_B} 是 \texttt{1}。 範例
\thechapter.4 模擬一場只有兩位候選人的公投：讀入兩個得票數，決定勝者，
再印出勝者代號（\texttt{A} 或 \texttt{B}）與總票數。 邏輯跟「比大小」
沒什麼兩樣，差別在於 \texttt{winner} 的型別是 \texttt{Candidate} 而
不是裸的 \texttt{int}，我們比較時用的是 \texttt{CAND\_A} 這種有名字的
常數，而不是硬記「0 代表 A、1 代表 B」。

\begin{example}{c}{(referendum-winner.c) 用 enum 表示候選人並決定勝者}
#include <stdio.h>

typedef enum { CAND_A, CAND_B } Candidate;

int main(void) {
    int votesA, votesB;
    scanf("%d %d", &votesA, &votesB);

    Candidate winner;
    if (votesA > votesB) {
        winner = CAND_A;
    } else {
        winner = CAND_B;
    }

    char label = (winner == CAND_A) ? 'A' : 'B';
    printf("%c %d\n", label, votesA + votesB);
}
\end{example}

\iosample{
100 200 \\
}{
B 300 \\
}

以範例的輸入為例：\texttt{votesA = 100}、\texttt{votesB = 200}，條件
\texttt{votesA > votesB} 不成立，所以 \texttt{winner = CAND\_B}，對應
的字元代號是 \texttt{'B'}，總票數 \texttt{100 + 200 = 300}，輸出
\texttt{B 300}。 如果改成 \texttt{500 123}，則 \texttt{winner} 會是
\texttt{CAND\_A}，輸出 \texttt{A 623}。

\stylepoint{需要把 enum 常數轉成字元或字串給 \texttt{printf} 時，通常
還得另外寫一個 \texttt{if} 或 \texttt{switch}；enum 解決的是「程式內部用
具名常數做判斷」，不是自動變成可印的字串。}

\begin{drills}[label=drills:ext-enum]
\item \texttt{typedef enum \{ CAND\_A, CAND\_B \} Candidate;} 沒有特別
指定值時，\texttt{CAND\_A} 和 \texttt{CAND\_B} 分別是多少？
\item 範例讀入 \texttt{500 123} 時，\texttt{winner} 會是哪個常數？ 輸出
是什麼？
\item 能不能直接寫 \texttt{printf("\%s", winner)} 印出 \texttt{"CAND\_A"}
這樣的名字？ 為什麼？
\end{drills}

\section{union：同一段記憶體的不同解讀}
\label{sec:ext-union}

\texttt{struct} 把多個欄位\emph{各自}占一段記憶體排在一起；\texttt{union}
則讓多個欄位\emph{共用}同一段記憶體。 同一時刻只有其中一個欄位「有效」，
但你可以用不同的型別去\emph{解讀}同一段位元組。 這在 low-level 程式裡很
常見：例如把一個 \texttt{int} 拆成四個 \texttt{unsigned char} 去看每一
個位元組，或把浮點數的位元模式當整數讀出來（這在浮點數那一章我們刻意
沒做，因為牽涉 IEEE 754 細節，但 union 正是標準做法）。

\begin{code}{c}{union 讓 int 與 char 陣列共用記憶體}
union IntBytes {
    int value;
    unsigned char bytes[4];
};

union IntBytes x;
scanf("%d", &x.value);
printf("%u\n", x.bytes[0]);
\end{code}

union 的大小，等於它所有成員中\emph{最大}的那一個所占的空間（這裡是
\texttt{int} 的 $4$ 位元組）。 寫入 \texttt{x.value} 之後，再讀
\texttt{x.bytes[0]} 到 \texttt{x.bytes[3]}，看到的其實是同一個
\texttt{int} 在記憶體裡的逐位元組分解。 順序取決於 CPU 是 little-endian
還是 big-endian，這也是為什麼 union 常跟「記憶體布局」一起教。

\begin{center}
\begin{tikzpicture}
    \node[hlcell] (b0) {b0};
    \node[hlcell, right=0pt of b0] (b1) {b1};
    \node[hlcell, right=0pt of b1] (b2) {b2};
    \node[hlcell, right=0pt of b2] (b3) {b3};
    \coordinate (bMid) at ($(b0.south)!0.5!(b3.south)$);
    \node[idxlbl, below=0.2cm of bMid] {\texttt{x.value}（同一個 int）};
    \node[idxlbl, anchor=west, align=left, right=0.35cm of b3] {也可寫成\\ \texttt{x.bytes[0..3]}};
\end{tikzpicture}
\diagcaption{\texttt{union IntBytes} 的四個位元組與單一 \texttt{int}
共用同一段記憶體；改 \texttt{value} 會連帶改變 \texttt{bytes} 的內容。}
\label{fig:ch21-union-bytes}
\end{center}

範例 \thechapter.6 讀入一個整數，依序印出它的四個位元組（由低位元組到
高位元組）。 在常見的 little-endian 機器上，輸入 \texttt{256} 會得到
\texttt{0 1 0 0}：因為 $256 = 0 \times 1 + 1 \times 256 + 0 + 0$，最低
位元組是 \texttt{0}，次低位元組是 \texttt{1}。

\begin{example}{c}{(union-bytes.c) 用 union 檢視 int 的各個位元組}
#include <stdio.h>

union IntBytes {
    int value;
    unsigned char bytes[4];
};

int main(void) {
    union IntBytes x;
    scanf("%d", &x.value);

    for (int i = 0; i < 4; i++) {
        if (i) {
            printf(" ");
        }
        printf("%u", x.bytes[i]);
    }
    printf("\n");
}
\end{example}

\iosample{
256 \\
}{
0 1 0 0 \\
}

\keypoint{union 的所有成員共享同一段記憶體；寫入一個成員後再讀另一個成
員，讀到的是同一位元模式的\emph{不同解讀}，不是複製出來的第二份資料。}

\begin{drills}[label=drills:ext-union]
\item \texttt{union IntBytes \{ int value; unsigned char bytes[4]; \}}
這個 union 的 \texttt{sizeof} 是多少？
\item 在 little-endian 機器上，輸入 \texttt{256} 給 \texttt{x.value}，
\texttt{x.bytes[0]} 到 \texttt{x.bytes[3]} 分別是多少？
\item \texttt{struct} 跟 \texttt{union} 的差別是什麼——各自的成員各占一
段記憶體，還是共用同一段？
\end{drills}

\section{storage class：變數能活多久、誰能看到}
\label{sec:ext-storage-class}

每個變數除了型別之外，還有一個 storage class（儲存類別），決定它「在
哪裡分配」「能活多久」「哪些函式看得到」。 初學時最常碰到的是
\texttt{auto}（區域變數的預設）跟 \texttt{static}；之後寫多檔案程式會再
遇到 \texttt{extern}。

\begin{code}{c}{區域 auto 變數 vs. static 區域變數}
void foo() {
    int autoVar = 0;      // 每次呼叫都重新建立，離開函式就消失
    autoVar++;

    static int persist = 0;  // 只初始化一次，值在多次呼叫之間保留
    persist++;
}
\end{code}

\texttt{auto} 區域變數存在每次函式呼叫的 stack frame 裡：函式一返回，這
一層堆疊就被整個丟棄，變數也跟著消失。 \texttt{static} 區域變數雖然
「寫在函式裡面」，但記憶體不在 stack 上反覆建立／銷毀。 編譯器只會在
程式啟動時初始化一次，之後每次呼叫都沿用同一塊儲存空間。 這跟函式那一
章講的 call stack 可以對照：一般區域變數是跟著 stack frame 來來去去，
\texttt{static} 區域變數則像躲在函式後面的「永久抽屜」。

範例 \thechapter.8 連續呼叫 \texttt{countCalls} 三次，每次印出累計次
數 \texttt{1}、\texttt{2}、\texttt{3}。 如果把 \texttt{static int count}
改成一般的 \texttt{int count}，每次呼叫都會從 \texttt{0} 重新開始，三
次都會印 \texttt{1}。

\begin{example}{c}{(static-counter.c) static 區域變數保留呼叫次數}
#include <stdio.h>

void countCalls(void) {
    static int count = 0;
    count++;
    printf("%d\n", count);
}

int main(void) {
    int n;
    scanf("%d", &n);
    for (int i = 0; i < n; i++) {
        countCalls();
    }
}
\end{example}

\iosample{
3 \\
}{
1 \\
2 \\
3 \\
}

\texttt{extern} 則是「這個變數的\emph{定義}在別的檔案，這裡只是宣告我
要用它」。 片語 \thechapter.9 示範典型寫法：在標頭檔用 \texttt{extern}
宣告，在某個 \texttt{.c} 檔真正定義並初始化。 這讓多個原始檔可以共享
同一個全域變數，又不會在連結時因為重複定義而報錯。

\begin{code}{c}{extern 宣告 vs. 另一個檔案中的定義}
// shared.h
extern int globalCounter;

// fileA.c
#include "shared.h"
int globalCounter = 0;   // 唯一的定義

// fileB.c
#include "shared.h"
void bump() { globalCounter++; }  // 使用 extern 宣告的全域變數
\end{code}

\stylepoint{\texttt{static} 也可以修飾「檔案層級的全域變數或函式」，意
思變成「只在這個 \texttt{.c} 檔內可見」。 跟函式內的 \texttt{static}
區域變數是不同用法，但關鍵字相同，讀程式時要看它修飾的是區域變數還是
全域名稱。}

\begin{drills}[label=drills:ext-storage-class]
\item \texttt{countCalls} 裡的 \texttt{static int count} 若改成一般的
\texttt{int count}，連續呼叫三次會印出什麼？ 跟原本的 \texttt{1 2 3}
差在哪裡？
\item \texttt{auto} 區域變數存在哪裡？ 函式返回後它會發生什麼事？
\item \texttt{extern} 宣告和它真正的定義，通常分別放在什麼樣的檔案裡？
\end{drills}

\section{動態記憶體管理簡介}
\label{sec:ext-dynamic-memory}

到此為止，我們的變數都是在\textbf{編譯時}決定大小、在\textbf{執行時}建立
和銷毀。 但有時候程式需要「根據使用者輸入決定需要多少記憶體」，或「在
執行期間新增、刪除資料元素」。 這時就需要\textbf{動態配置}（dynamic
allocation）——在程式執行時，用 \texttt{malloc} 從 \emph{heap} 上要求一
塊記憶體，用完後用 \texttt{free} 歸還。 指標是動態配置的核心應用：沒
有指標，就無法把記憶體中四散各地、執行期間才建立的區塊連接起來，這正
是鏈結串列、樹這類\emph{資料結構}的基礎。

\texttt{malloc}／\texttt{free} 的完整規則——三種會壞掉的方式（記憶體
洩漏、懸置指標、雙重釋放）、怎麼用 AddressSanitizer 之類的工具把它們
抓出來、\texttt{calloc}／\texttt{realloc} 怎麼用——已經在
第~\ref{ch:datastructure}~章（資料結構）系統講過，這裡不重複；還沒讀
過的話，建議先讀那一章再回來看後面的案例。

\begin{drills}[label=drills:ext-dynamic-memory]
\item \texttt{malloc}／\texttt{free} 的完整規則（記憶體洩漏、懸置指
標、雙重釋放三種錯誤，以及用 AddressSanitizer 之類的工具找出它們）寫
在哪一章？
\item 為什麼說「指標是動態配置的核心應用」？
\end{drills}

\section{struct 的記憶體布局與對齊}
\label{sec:ext-struct-padding}

到這裡，我們已經知道 \texttt{int} 占 $4$ 位元組、\texttt{char} 占
$1$ 位元組。 但把不同型別的欄位\emph{放進同一個 struct} 時，總大小
往往\emph{不是}各欄位大小單純相加。 編譯器會在欄位之間插入 padding
（填充位元組），讓每個欄位的起始位址符合 CPU 的\emph{對齊}
（alignment）要求。 對齊不是 C 語言隨意加的累贅：許多 CPU 從「不是
$4$ 的倍數」的位址讀取 \texttt{int} 會比較慢，有些架構甚至直接當
成硬體錯誤。

先看一個只有三個 \texttt{char} 的 struct，沒有 padding 的餘地：

\begin{code}{c}{連續三個 char，大小就是 3}
struct ThreeChars {
    char a, b, c;
};
// sizeof(struct ThreeChars) == 3
\end{code}

如果在兩個 \texttt{char} 中間插入一個 \texttt{int}，布局就會變複雜：

\begin{center}
\begin{tikzpicture}
    \node[hlcell] (structA) {a};
    \node[cell, minimum width=0.45cm, right=0pt of structA] (pad1a) {};
    \node[cell, minimum width=0.45cm, right=0pt of pad1a] (pad1b) {};
    \node[cell, minimum width=0.45cm, right=0pt of pad1b] (pad1c) {};
    \node[hlcell2, right=0pt of pad1c] (x1) {x};
    \node[hlcell2, right=0pt of x1] (x2) {x};
    \node[hlcell2, right=0pt of x2] (x3) {x};
    \node[hlcell2, right=0pt of x3] (x4) {x};
    \node[hlcell, right=0pt of x4] (structC) {c};
    \node[cell, minimum width=0.45cm, right=0pt of structC] (pad2a) {};
    \node[cell, minimum width=0.45cm, right=0pt of pad2a] (pad2b) {};
    \node[cell, minimum width=0.45cm, right=0pt of pad2b] (pad2c) {};

    \node[idxlbl, below=0.2cm of structA] {a};
    \node[idxlbl, below=0.2cm of pad1b] {pad};
    \coordinate (xMid) at ($(x1.south)!0.5!(x4.south)$);
    \node[idxlbl, below=0.2cm of xMid] {x (int)};
    \node[idxlbl, below=0.2cm of structC] {c};
    \node[idxlbl, below=0.2cm of pad2b] {pad};
\end{tikzpicture}
\diagcaption{\texttt{struct CharIntChar \{ char a; int x; char c; \}} 在
常見 $32$-bit 平台上的布局：\texttt{int x} 必須從 $4$ 的倍數位址開始，
所以 \texttt{a} 後面插入 $3$ 個 padding 位元組，末尾也可能再 pad 到
整個 struct 大小是 $4$ 的倍數。}
\label{fig:ch21-struct-padding}
\end{center}

範例 \thechapter.11 用 \texttt{sizeof} 直接量出兩個 struct 的大小：
\texttt{struct ThreeChars} 是 \texttt{3}，\texttt{struct CharIntChar} 在
同一台機器上是 \texttt{12}。 中間那個 \texttt{int} 加上前後 padding，
把三個「邏輯上只有 $1 + 4 + 1 = 6$ 位元組」的欄位撐成 $12$ 位元組。

\begin{example}{c}{(struct-padding.c) 用 sizeof 觀察 padding 造成的差異}
#include <stdio.h>

struct ThreeChars {
    char a, b, c;
};

struct CharIntChar {
    char a;
    int x;
    char c;
};

int main(void) {
    printf("%zu %zu\n", sizeof(struct ThreeChars), sizeof(struct CharIntChar));
}
\end{example}

\iosample{}{
3 12 \\
}

\keypoint{\texttt{sizeof(struct ...)} 量的是「編譯器實際分配的大小」，包
含 padding；不能假設 struct 大小等於各欄位大小之和。 這在之後用
\texttt{memcpy} 複製 raw 位元組、或跟硬體暫存器對接時特別重要。}

如果真的很在意每個位元組都不能浪費（例如嵌入式協定封包），可以請編譯
器用 \texttt{\_\_attribute\_\_((packed))} 或 \texttt{\#pragma pack}
關閉 padding。 但這可能導致未對齊存取變慢，甚至觸發硬體例外，所以一般
應用程式很少這樣做，比較常見的是\emph{調整欄位順序}（把相同型別或
較大的欄位放一起）來減少 padding，而不是完全關掉對齊。

\begin{drills}[label=drills:ext-struct-padding]
\item \texttt{struct ThreeChars \{ char a, b, c; \}} 的 \texttt{sizeof}
是多少？ 為什麼沒有 padding？
\item \texttt{struct CharIntChar \{ char a; int x; char c; \}} 在範例機
器上 \texttt{sizeof} 是 \texttt{12}，不是邏輯上的 \texttt{6}，多出來的
$6$ 個位元組是什麼？
\item 為什麼 \texttt{int} 欄位要從 $4$ 的倍數位址開始，編譯器才要插入
padding？
\end{drills}

\section{整數在記憶體裡長什麼樣}
\label{sec:ext-int-bytes}

把前面幾節串起來：一個 \texttt{int} 在記憶體裡就是連續的 $4$ 個位元
組；union 讓我們用 \texttt{bytes[0..3]} 逐一檢視；struct padding 則告訴
我們「型別大小」跟「布局規則」會一起決定位元組怎麼排。 下一章的位元
遮罩、移位、popcount，全都建立在「整數就是一串固定長度的位元，從最低
位（LSB）到最高位（MSB）有明確編號」這個圖像上。 如果這裡的 union 範
例 \thechapter.6 還記得 \texttt{256} 拆成 \texttt{0 1 0 0}，下一章讀
\texttt{x \& 1}、\texttt{x >> 1} 時就不會覺得那些運算子從天而降。

\begin{center}
\begin{tikzpicture}
    \foreach \v/\lbl [count=\i from 0] in {0/b0 (LSB), 1/b1, 0/b2, 0/b3 (MSB)} {
        \ifnum\i=0
            \node[cell] (byte\i) {\v};
        \else
            \pgfmathtruncatemacro{\prev}{\i-1}
            \node[cell, right=0pt of byte\prev] (byte\i) {\v};
        \fi
        \node[idxlbl, below=0.2cm of byte\i] {\lbl};
    }
    \node[idxlbl, right=0.65cm of byte3] {值 $= 256$};
\end{tikzpicture}
\diagcaption{little-endian 機器上 \texttt{int} 值 \texttt{256} 的四個位
元組；位元運算章節會改從「位元 0、位元 1、\ldots」的角度看同一串
資料。}
\label{fig:ch21-little-endian-bytes}
\end{center}

\begin{drills}[label=drills:ext-int-bytes]
\item 這一節說「整數就是一串固定長度的位元，從 LSB 到 MSB 有明確編
號」，這個圖像跟哪一節的 union 範例直接對應？
\item little-endian 機器上，\texttt{256} 的四個位元組由低到高分別是
什麼？
\end{drills}

\section[位元操作（見第 12 章）]{位元操作詳解}
\label{sec:ext-bitops-pointer}

位元操作的完整內容在獨立的第~\ref{ch:bitoperations}~章（位元操作）：
進制轉換、操作符定義，到打包、解包與演算法應用都在那裡。 本章只在後
面的案例裡把位元操作跟其他高等型別結合，示範實戰應用，不重複講一次
基礎，詳見第~\ref{ch:bitoperations}~章。

\begin{drills}[label=drills:ext-bitops-pointer]
\item 位元操作的完整內容（進制轉換、操作符定義、打包解包與演算法應
用）在哪一章？
\item 這一節為什麼不重新講一次位元操作的基礎？
\end{drills}

\section{高等技術的綜合應用}
\label{sec:ext-synthesis}

前面幾節的觀念——型別、記憶體布局、位元運算、動態配置——在實務中往往
一起用。 這一節用三個具體例子展示「高等技術」怎麼組織出有用的東西。

\subsection{案例 1：集合型別（bitset）}

題庫裡常見「集合」題目：給定元素上限（例如最多 64 個元素），用位元表示
某元素是否在集合裡。 一個 \texttt{unsigned long long} 可以表示 64 個
元素的集合。

\begin{code}{c}{集合型別的設計}
typedef unsigned long long Set;

// 包含元素 i
Set insert(Set s, int i) {
    return s | (1ULL << i);
}

// 移除元素 i（不叫 remove：那是 <stdio.h> 刪檔案函式的名字，撞名要避免）
Set set_remove(Set s, int i) {
    return s & ~(1ULL << i);
}

// 檢查元素 i 是否存在
int contains(Set s, int i) {
    return (s >> i) & 1;
}

// 聯集
Set union_set(Set a, Set b) {
    return a | b;
}

// 交集
Set intersect(Set a, Set b) {
    return a & b;
}
\end{code}

這個設計結合了 \texttt{typedef}（提供有語意的名字）、位元運算（高效
實作）、以及位元字段的思想（每位代表一個元素）。

\subsection{案例 2：結構序列化}

讀寫二進位檔案或網路協定時，常常要把 \texttt{struct} 的位元模式逐位元
組寫出去。 這時會結合 \texttt{union} 和 \texttt{memcpy}：

\begin{code}{c}{用 union 進行結構序列化}
struct Color {
    unsigned char r, g, b;
};

union ColorBytes {
    struct Color rgb;
    unsigned int packed;  // 32-bit integer
};

// 把 RGB 顏色轉成 32-bit 整數（little-endian）
unsigned int colorToInt(struct Color c) {
    union ColorBytes cb;
    cb.rgb = c;
    return cb.packed;
}
\end{code}

這裡結合了 \texttt{struct} 的語意清晰度、\texttt{union} 的位元重新解讀，
以及 \texttt{sizeof} 對記憶體布局的理解。

\subsection{案例 3：鏈結串列節點}

鏈結串列是資料結構的基礎，每個節點是用 \texttt{malloc} 動態建立的
\texttt{struct}，內含指向下一節點的指標。 這同時用上了 \texttt{typedef}、
\texttt{struct}、指標，以及動態配置：

\begin{code}{c}{鏈結串列節點的典型設計}
typedef struct Node {
    int data;
    struct Node *next;  // 自指型別
} Node;

Node *head = NULL;

// 新增節點到鏈結串列前端
void insert(int val) {
    Node *new_node = (Node *)malloc(sizeof(Node));
    if (new_node == NULL) return;
    new_node->data = val;
    new_node->next = head;
    head = new_node;
}

// 走訪鏈結串列
void print_all(void) {
    for (Node *p = head; p != NULL; p = p->next) {
        printf("%d ", p->data);
    }
}
\end{code}

這三個案例展示了高等技術不是為了複雜而複雜，而是精確地用編譯器和 CPU
的能力去解決問題：集合用位元的高效性，序列化用 union 的位元重新解讀，
鏈結串列用指標和動態配置的靈活性。 掌握這些工具，就能寫出既高效又清晰
的程式碼。

\begin{drills}[label=drills:ext-synthesis]
\item 案例 1 的 bitset 為什麼把「移除元素」的函式命名為
\texttt{set\_remove} 而不是直接叫 \texttt{remove}？
\item 案例 2 的 \texttt{union ColorBytes} 讓 \texttt{struct Color} 跟
\texttt{unsigned int packed} 共用記憶體，這裡 union 的用途，跟前面
\texttt{union IntBytes} 那節的用途是同一件事嗎？
\item 案例 3 的 \texttt{Node} 結構裡，\texttt{struct Node *next;} 為什
麼可以在自己的定義裡出現指向自己型別的指標，卻不會造成結構大小是無
窮大？
\end{drills}

\addcontentsline{toc}{section}{練習解答}
\section*{練習解答}

各節練習的答案收在這裡。 短題只給能一句話核對的結果。

\paragraph{第~\ref{sec:ext-typedef}~節}
\begin{enumerate}
\item 底層型別仍是 \texttt{int}；編譯器把 \texttt{Score} 當成
\texttt{int} 處理，不會多占記憶體或改變運算規則。
\item 只需要改 \texttt{typedef} 那一行本身；這正是 typedef 讓「物理表
示可以換、語意名字留著」的價值。
\item 一樣，兩者都是替既有型別（一個是 \texttt{int}、一個是
\texttt{struct}）取一個表達語意的別名。
\end{enumerate}

\paragraph{第~\ref{sec:ext-enum}~節}
\begin{enumerate}
\item \texttt{CAND\_A} 是 \texttt{0}，\texttt{CAND\_B} 是 \texttt{1}。
\item \texttt{winner} 是 \texttt{CAND\_A}，輸出 \texttt{A 623}。
\item 不行；enum 常數只是有名字的整數，本身不會自動變成可印的字串，
要另外寫 \texttt{if} 或 \texttt{switch} 轉換。
\end{enumerate}

\paragraph{第~\ref{sec:ext-union}~節}
\begin{enumerate}
\item 等於最大成員 \texttt{int value} 的大小，也就是 \texttt{4}。
\item \texttt{0 1 0 0}（\texttt{256} 的最低位元組是 \texttt{0}，次低位
元組是 \texttt{1}）。
\item \texttt{struct} 的成員各自占一段記憶體；\texttt{union} 的成員共
用同一段記憶體，同一時刻只有其中一個有效。
\end{enumerate}

\paragraph{第~\ref{sec:ext-storage-class}~節}
\begin{enumerate}
\item 會連續印出 \texttt{1 1 1}，因為每次呼叫都重新從 \texttt{0} 開
始；原本的 \texttt{static} 版本則保留累計值，印出 \texttt{1 2 3}。
\item 存在函式呼叫的 stack frame 裡；函式返回後這層堆疊被丟棄，變數
跟著消失。
\item \texttt{extern} 宣告通常放在標頭檔（\texttt{.h}），真正的定義放
在某一個 \texttt{.c} 檔裡。
\end{enumerate}

\paragraph{第~\ref{sec:ext-dynamic-memory}~節}
\begin{enumerate}
\item 第~\ref{ch:datastructure}~章（資料結構）。
\item 因為執行期間才建立、分散在記憶體各處的區塊，必須靠指標才能連
接起來，這正是鏈結串列、樹這類資料結構的基礎。
\end{enumerate}

\paragraph{第~\ref{sec:ext-struct-padding}~節}
\begin{enumerate}
\item \texttt{3}；三個 \texttt{char} 都是 $1$ 位元組，彼此對齊要求相
同，不需要 padding。
\item padding：\texttt{a} 後面補 $3$ 個位元組讓 \texttt{int x} 對齊到
$4$ 的倍數位址，結尾可能再補 $3$ 個位元組讓整個 struct 大小是 $4$ 的
倍數。
\item 因為許多 CPU 從不是 $4$ 的倍數的位址讀取 \texttt{int} 會比較慢，
甚至可能直接觸發硬體錯誤，編譯器插入 padding 讓每個欄位起始位址符合
對齊要求。
\end{enumerate}

\paragraph{第~\ref{sec:ext-int-bytes}~節}
\begin{enumerate}
\item 跟 \texttt{union IntBytes} 那節、把 \texttt{256} 拆成
\texttt{0 1 0 0} 四個位元組的範例對應。
\item \texttt{0}、\texttt{1}、\texttt{0}、\texttt{0}（由低到高）。
\end{enumerate}

\paragraph{第~\ref{sec:ext-bitops-pointer}~節}
\begin{enumerate}
\item 第~\ref{ch:bitoperations}~章（位元操作）。
\item 因為完整的基礎已經在第~\ref{ch:bitoperations}~章系統講過，這一
節只在案例裡把位元操作跟其他高等型別結合示範應用，重複講一次沒有必
要。
\end{enumerate}

\paragraph{第~\ref{sec:ext-synthesis}~節}
\begin{enumerate}
\item 因為 \texttt{remove} 是 \texttt{<stdio.h>} 裡刪除檔案的函式名
字，跟它撞名容易造成混淆或連結衝突，所以改用 \texttt{set\_remove}。
\item 是同一件事：兩處的 union 都是讓一段記憶體同時用兩種型別解讀，
藉此把資料在不同表示法之間轉換或拆解。
\item 不會，因為 \texttt{struct Node} 裡放的是指向 \texttt{struct Node}
的\emph{指標}，指標本身大小固定（不隨指向的型別而變），所以不會造成
無窮遞迴展開。
\end{enumerate}
